\subsection{Continuous visibility correction}\label{sec:autovis}

Water absorbs red light first, scatters light into the line of sight, and attenuates
contrast with distance \cite{jaffe1990,schettini2010}: underwater pictures are
blue-green, compressed in brightness and hazy, and they change as the vehicle
descends. Because the pilot cannot tune effects while flying, \eeye{} provides a
correction, off by default and switched on from a single control on the page, that
adapts continuously and has no parameters to set. It runs before the operator chain,
on the pilot's picture only; \emph{raw} is never altered, so that autonomy code and the
pilot's own comparison see the camera's frame.

The correction is steered by statistics of the raw frame, taken on every second pixel
of every second row: the luma histogram, and the mean chroma offsets $\bar u$,
$\bar v$ from neutral over samples with $24 \leq Y \leq 235$, which excludes pixels
too dark to carry colour and pixels that have lost it to clipping. From these it
derives three corrections, each a target that the applied correction follows with the
first-order filter
\begin{equation}\label{eq:smooth}
  c_k = c_{k-1} + \alpha_k\,(c^{*}_k - c_{k-1}), \qquad
  \alpha_k = 1 - e^{-\Delta t_k/\tau},
\end{equation}
where $\Delta t_k$ is the difference between the capture times of consecutive frames,
so that the response does not depend on the frame rate. A gap of more than 1~s, or
switching the correction on, sets $c_k = c^{*}_k$, without fading in from no correction.

\emph{Colour cast.} Following the grey-world assumption \cite{buchsbaum1980}, the
target chroma shift removes the measured cast down to a deadband $d$,
\begin{equation}\label{eq:cast}
  (\Delta u^{*}, \Delta v^{*}) = -\frac{\max(m - d, 0)}{m}\,(\bar u, \bar v),
  \qquad m = \sqrt{\bar u^2 + \bar v^2},
\end{equation}
with each component limited to $\pm 48$. The deadband leaves the small cast that a
scene's own colours produce; without it the synthetic scene of Figure~\ref{fig:autovis}
(red, yellow and orange objects, cast 9.9) had its grey background shifted by 11.7
chroma units towards blue. Since a cast is a gain on the light, the shift applied at a
chroma sample is scaled by $\min(Y/\bar Y, 1.5)$, where $Y$ is the luma under the
sample and $\bar Y$ the mean luma of the measured samples, and chroma is then scaled
about 128 by a saturation gain of 1.2.

\emph{Brightness range.} The 0.5th and 99.5th luma percentiles, $p_\mathrm{lo}$ and
$p_\mathrm{hi}$, are stretched to 0 and 255; percentiles rather than extrema, so that
a few bright particles of backscatter cannot defeat the stretch. When
$p_\mathrm{hi} - p_\mathrm{lo}$ is under 96 levels the interval is widened to 96 about
its centre, which limits the gain to 2.7 and keeps the sensor noise of a flat, dark
frame from being amplified into texture. A gamma $\gamma = \ln 0.45 / \ln \tilde m$,
limited to 0.55--1.6, where $\tilde m$ is the median luma after the stretch, places
the median at 45\% of full scale.

\emph{Haze.} Contrast-limited adaptive histogram equalization (CLAHE)
\cite{pizer1987,zuiderveld1994}, applied to luma as in earlier underwater
enhancement \cite{hitam2013}, restores local contrast. The frame is divided into
tiles of about 160$\times$120 pixels (at most 8$\times$6); each tile's histogram is
clipped at twice the height of a flat histogram, the excess is spread evenly over all
levels, and the cumulative histogram gives the tile's mapping, mixed with the identity
at a weight of 0.7. The tile mappings are smoothed over time by Eq.~\eqref{eq:smooth},
composed with the global stretch, and interpolated bilinearly between the four nearest
tile centres for each pixel. The tile histograms are taken from the source luma and
pushed through the global mapping, so that the globally corrected image need not be
written and read again. Table~\ref{tab:autovis} lists the constants.

The correction was evaluated on synthetic frames formed by the classical underwater
image model \cite{jaffe1990}, a simplification of the revised model of
Akkaynak and Treibitz \cite{akkaynak2018},
\begin{equation}\label{eq:water}
  I_c = J_c\,t_c + B_c\,(1 - t_c), \qquad c \in \{R, G, B\},
\end{equation}
with transmissions $t = (0.25, 0.70, 0.80)$ and veiling light $B = (15, 110, 125)$,
applied to a scene of coloured patches, gratings of three spatial frequencies, a
gradient and noise of $\pm 2$ levels (Figure~\ref{fig:autovis}).
Section~\ref{sec:res-autovis} gives the results.

\begin{table}[!t]
  \centering\small
  \caption{Constants of the visibility correction. They were set on the synthetic
    frames of Eq.~\eqref{eq:water} and a laptop camera, and are to be revised on
    footage from the vehicle.}\label{tab:autovis}
  \begin{tabular}{@{}L{5.4cm}r@{}}
    \toprule
    Constant & Value \\
    \midrule
    Smoothing time constant $\tau$, Eq.~\eqref{eq:smooth} & 0.4~s \\
    Gap that restarts the smoothing & 1~s \\
    Luma range of samples used for the cast & 24--235 \\
    Cast deadband $d$, Eq.~\eqref{eq:cast} & 8 \\
    Largest chroma shift & $\pm 48$ \\
    Saturation gain & 1.2 \\
    Stretch percentiles & 0.5th, 99.5th \\
    Smallest stretched interval (largest gain) & 96 levels (2.7) \\
    Target median & 0.45 \\
    Gamma limits & 0.55--1.6 \\
    CLAHE tile size (largest grid) & $\approx$160$\times$120~px (8$\times$6) \\
    CLAHE clip limit & 2.0 \\
    Weight of the local mapping & 0.7 \\
    \bottomrule
  \end{tabular}
\end{table}
